% ===========================================================================
\section{反例：不是 UFD 的整环}\label{sec:counter}

\begin{definition}[{环的生成 $R[u]$}]
设 $R\subseteq R'$ 为环，$u\in R'$。$R[u]$ 表示 $R'$ 中包含 $R$ 与 $u$
的\keyword{最小子环}：
\[
  R[u] = \{\, r_0 + r_1 u + \cdots + r_t u^t : r_i\in R,\ t\ge 0\,\}.
\]
\end{definition}

于是 $\Z[\sqrt{-5}] = \{\,a+b\sqrt{-5} : a,b\in\Z\,\}$，
$\Z[\sqrt{10}] = \{\,a+b\sqrt{10} : a,b\in\Z\,\}$。它们都是整环：交换性
显然；无零因子因为它们是域（$\C$、$\R$）的子环。

\subsection{\texorpdfstring{$\Z[\sqrt{-5}]$}{Z[√-5]} 中分解不唯一}

\begin{example}[{$\Z[\sqrt{-5}]$ 中分解不唯一}]
\[
  6 = 2\times 3 = (1+\sqrt{-5})(1-\sqrt{-5}).
\]
各因子不可约的验证利用几何图像：把 $\Z[\sqrt{-5}]$ 看成复平面上的整点格。
模长满足 $|xy| = |x||y|$；离 $0$ 最近的非零格点是 $\pm 1$，故单位只有
$\pm 1$。找一个元素的因子只需在模长更小的格点中找——以原点为圆心、
$|a|$ 为半径的圆内只有有限个格点，逐一枚举即可验证 $2, 3,
1\pm\sqrt{-5}$ 均不可约。而 $2$ 与 $1\pm\sqrt{-5}$ 互不相伴（单位只有
$\pm 1$），两个分解之间没有相伴对应，故 $\Z[\sqrt{-5}]$ 不是 UFD。
\end{example}

这同时给出「不可约但不是素元」的例子：$2\mid 6 =
(1+\sqrt{-5})(1-\sqrt{-5})$，但 $2\nmid 1\pm\sqrt{-5}$。

\subsection{\texorpdfstring{$\Z[\sqrt{10}]$}{Z[√10]} 中分解不唯一}

\begin{example}[{$\Z[\sqrt{10}]$ 中分解不唯一}]
\[
  6 = 2\times 3 = (4+\sqrt{10})(4-\sqrt{10})\quad (16-10=6),
\]
\[
  15 = 3\times 5 = (5+\sqrt{10})(5-\sqrt{10})\quad (25-10=15).
\]
困难在于：$\{a+b\sqrt{10}\}$ 在实数轴上\keyword{稠密}，普通绝对值下
「模长更小的元素」有无穷多个，几何枚举法失效。解决办法是重新定义
「大小」——\keyword{范数（Norm）}：
\[
  N(a+b\sqrt{10}) = |a^2 - 10b^2|.
\]
范数是积性的，$N(x)N(y) = N(xy)$，可直接展开验证：
\[
  (a^2-10b^2)(c^2-10d^2) = (ac+10bd)^2 - 10\,(ad+bc)^2.
\]
因此 $x\mid y \Rightarrow N(x)\mid N(y)$。于是任何元素的潜在因子只有
范数严格更小的那些；范数取值为自然数，给定范数值的候选只有有限多种
（对 $a^2 - 10b^2$ 的给定值枚举 $a,b$）。由此同样可验证上面的因子均
不可约、两种分解互不相伴，故 $\Z[\sqrt{10}]$ 不是 UFD。
\end{example}
